\documentclass[11pt,a4paper]{article}

%=======================================================================
% Préambule
%=======================================================================

\usepackage[T1]{fontenc}
\usepackage[utf8]{inputenc}
% natbib est chargé avant babel-french, qui l'exige pour adapter ses
% redéfinitions de la bibliographie.
\usepackage[round]{natbib}
% La bibliographie constitue l'annexe D : son titre est produit comme une
% section numérotée plutôt que comme la section étoilée par défaut.
\renewcommand{\bibsection}{\section{Bibliographie}\label{app:biblio}\medskip}
\usepackage[french]{babel}
\usepackage{lmodern}
\usepackage{microtype}
\usepackage[margin=2.4cm]{geometry}
\usepackage{amsmath,amssymb,amsthm,mathtools}
\usepackage{booktabs}
\usepackage{tabularx}
% Colonne extensible non justifiée : les colonnes étroites des tableaux
% comparatifs produiraient sinon des espacements inter-mots très irréguliers.
\newcolumntype{Y}{>{\raggedright\arraybackslash}X}
\usepackage{enumitem}
\usepackage{xcolor}
\usepackage{tikz}
\usetikzlibrary{arrows.meta,positioning,shapes.geometric,calc,fit,backgrounds}
\usepackage{algorithm}
\usepackage{algpseudocode}
\usepackage{listings}

% Les couleurs sont définies AVANT le chargement de hyperref, qui les
% référence dans ses options ; hidelinks est incompatible avec colorlinks.
\definecolor{Bleu}{HTML}{1F4E79}
\definecolor{Gris}{HTML}{6B7280}
\definecolor{GrisClair}{HTML}{F1F2F4}
\definecolor{Rouge}{HTML}{A83232}
\definecolor{Vert}{HTML}{2E6B4F}

\usepackage[colorlinks=true,linkcolor=Bleu,citecolor=Bleu,urlcolor=Bleu]{hyperref}

% --- Environnements de théorèmes ---
\theoremstyle{definition}
\newtheorem{definition}{Définition}[section]
\newtheorem{propriete}{Propriété}[section]
\newtheorem{hypothese}{Hypothèse}[section]
\theoremstyle{remark}
\newtheorem{remarque}{Remarque}[section]
\newtheorem{avertissement}{Avertissement}[section]

% --- Algorithmes en français ---
\floatname{algorithm}{Algorithme}
\renewcommand{\listalgorithmname}{Liste des algorithmes}
\algrenewcommand\algorithmicrequire{\textbf{Entrées :}}
\algrenewcommand\algorithmicensure{\textbf{Sortie :}}
\algrenewcommand\algorithmicfor{\textbf{pour}}
\algrenewcommand\algorithmicforall{\textbf{pour tout}}
\algrenewcommand\algorithmicdo{\textbf{faire}}
\algrenewcommand\algorithmicwhile{\textbf{tant que}}
\algrenewcommand\algorithmicend{\textbf{fin}}
\algrenewcommand\algorithmicif{\textbf{si}}
\algrenewcommand\algorithmicthen{\textbf{alors}}
\algrenewcommand\algorithmicelse{\textbf{sinon}}
\algrenewcommand\algorithmicreturn{\textbf{retourner}}

% --- Listings Python ---
\lstset{
  language=Python,
  basicstyle=\ttfamily\footnotesize,
  keywordstyle=\color{Bleu}\bfseries,
  commentstyle=\color{Vert}\itshape,
  stringstyle=\color{Rouge},
  backgroundcolor=\color{GrisClair},
  frame=single,
  rulecolor=\color{Gris},
  breaklines=true,
  showstringspaces=false,
  columns=fullflexible,
  extendedchars=true,
  literate=%
    {é}{{\'e}}1 {è}{{\`e}}1 {ê}{{\^e}}1 {à}{{\`a}}1 {ù}{{\`u}}1
    {ç}{{\c c}}1 {î}{{\^i}}1 {ô}{{\^o}}1 {û}{{\^u}}1 {É}{{\'E}}1
    {â}{{\^a}}1 {ë}{{\"e}}1 {ï}{{\"i}}1 {°}{{$^\circ$}}1 {œ}{{\oe}}1
}

% --- Raccourcis mathématiques ---
\newcommand{\R}{\mathbb{R}}
\newcommand{\N}{\mathbb{N}}
\newcommand{\Z}{\mathbb{Z}}
\newcommand{\PP}{\mathbb{P}}
\newcommand{\EE}{\mathbb{E}}
\newcommand{\Unif}{\mathcal{U}}
\newcommand{\Bball}{\mathbb{B}}
\newcommand{\Sphere}{\mathbb{S}}
\newcommand{\norm}[1]{\left\lVert #1 \right\rVert}
\newcommand{\abs}[1]{\left\lvert #1 \right\rvert}
\newcommand{\scal}[2]{\left\langle #1 ,\, #2 \right\rangle}
\newcommand{\ind}[1]{\mathbf{1}\{#1\}}
\newcommand{\simplex}{\Delta^{d-1}}
\newcommand{\Tmap}{\mathbf{T}}
\newcommand{\pareto}{\succ_{\mathrm{P}}}
\newcommand{\paretoeq}{\succeq_{\mathrm{P}}}
\newcommand{\Xor}{\mathring{X}}
\newcommand{\med}{\operatorname{med}}
\newcommand{\MAD}{\operatorname{MAD}}
\newcommand{\rg}{\operatorname{rg}}

\title{\textbf{Méthodologie de synthèse multicritère}\\[0.3em]
\large Ordonner un ensemble de vecteurs de $\R^d$\\ sans pondération fixée \emph{a priori}}
\author{Note méthodologique}
\date{\today}

\begin{document}
\maketitle

\begin{abstract}
\noindent
Cette note décrit les méthodes permettant d'ordonner un ensemble fini de vecteurs de
$\R^d$ lorsque aucune hiérarchie n'est fixée \emph{a priori} entre les coordonnées,
c'est-à-dire lorsque tous les coefficients de la synthèse sont déterminés à partir des
données elles-mêmes. Six familles sont présentées : les ordres partiels issus de la
dominance, qui n'introduisent aucun coefficient ; les pondérations endogènes, qui
déduisent des poids d'une propriété statistique de l'échantillon ; les fonctions
d'agrégation, qui transforment un vecteur pondéré en un scalaire ; les rangs
multivariés fondés sur le transport optimal, qui ordonnent par la position dans le
nuage ; l'exploration de l'espace des poids, qui remplace un classement unique par une
distribution de classements ; enfin les procédures de comparaison et de contrôle, qui
mesurent l'accord entre méthodes et la stabilité des rangs. Chaque méthode est décrite
par son idée, ses hypothèses, ses définitions, son algorithme, ses propriétés
(monotonie, invariances, compensation, complexité), ses paramètres et ses limites. Le
dernier chapitre illustre l'ensemble sur un cas d'application.
\end{abstract}

\vspace{0.5em}
{\small\tableofcontents}
\newpage

%=======================================================================
\section{Introduction}
\label{sec:intro}
%=======================================================================

\medskip

\subsection{Le problème}

\medskip

On considère un ensemble fini de $n$ \emph{unités}, indexées par $i \in [n] := \{1,\dots,n\}$,
décrites chacune par $d$ \emph{métriques} indexées par $j \in [d]$. L'unité $i$ est
représentée par le vecteur ligne $x_i = (x_{i1},\dots,x_{id}) \in \R^d$ et l'ensemble des
observations forme la matrice

\medskip

\[
X = (x_{ij})_{i \in [n],\, j \in [d]} \in \R^{n \times d}.
\]

\bigskip

La question traitée est celle de la construction d'un ordre sur $\{x_1,\dots,x_n\}$ qui
rende compte simultanément des $d$ métriques, et dont tous les paramètres soient
\emph{estimés à partir de $X$} plutôt que fixés par jugement d'expert. Dans le cas le
plus courant, cet ordre est induit par une fonction de score $s : \R^d \to \R$, la
convention retenue ici étant qu'une valeur élevée de $s$ désigne une unité placée en tête
du classement. D'autres constructions ne produisent pas de score mais directement un ordre,
éventuellement partiel : l'ensemble des méthodes décrites dans cette note se répartit
entre ces deux sorties.

\bigskip

\subsection{L'absence d'ordre canonique sur \texorpdfstring{$\R^d$}{R\textasciicircum d}}

\medskip

Pour $d = 1$, l'ordre usuel de $\R$ répond à la question sans aucune convention
supplémentaire. Pour $d \ge 2$, une telle réponse n'existe pas : il n'y a pas d'ordre
total sur $\R^d$ qui soit à la fois compatible avec la structure d'espace vectoriel et
invariant par permutation des coordonnées.

\medskip

\begin{propriete}[Impossibilité d'un ordre total canonique]
  \label{prop:impossibilite}
  Soit $d \ge 2$ et soit $\preceq$ un ordre total sur $\R^d$ vérifiant les trois
  conditions suivantes :

  \medskip

  \begin{enumerate}[nosep,leftmargin=1.8em]
    \item \emph{compatibilité à l'addition} : $x \preceq y \Rightarrow x + z \preceq y + z$
    pour tout $z \in \R^d$ ;
    \item \emph{compatibilité à la multiplication scalaire positive} :
    $x \preceq y \Rightarrow \lambda x \preceq \lambda y$ pour tout $\lambda > 0$ ;
    \item \emph{invariance par permutation} : $x \preceq y \Rightarrow
    \sigma(x) \preceq \sigma(y)$ pour toute permutation $\sigma$ des coordonnées.
  \end{enumerate}

  \bigskip

  Alors $\preceq$ est l'ordre trivial, pour lequel tous les points sont équivalents.
\end{propriete}

\bigskip

En effet, notons $e_1,\dots,e_d$ la base canonique. L'invariance par permutation impose
que $e_1$ et $e_2$ soient équivalents, donc, par (1) appliqué avec $z = -e_2$, que
$e_1 - e_2$ soit équivalent à $0$. Par (3) à nouveau, $e_2 - e_1$ est également équivalent
à $0$, et par (1) tout multiple entier de $e_1 - e_2$ l'est aussi ; la combinaison de (1)
et (2) étend l'équivalence à la droite engendrée, puis, en itérant sur les autres paires
de coordonnées, à l'hyperplan $\{x : \sum_j x_j = 0\}$ tout entier. Un ordre total
compatible avec l'addition ne peut alors distinguer que la somme des coordonnées, laquelle
n'est pas invariante par changement d'échelle coordonnée par coordonnée ; imposer en outre
l'invariance d'échelle vide l'ordre de tout contenu.

\bigskip

La conséquence pratique est directe : tout classement complet des unités repose sur une
hypothèse supplémentaire, qui départage les situations où les métriques ne concordent pas.
Cette hypothèse est explicite lorsqu'un expert fixe des poids ; elle reste présente, sous
une autre forme, lorsque les poids sont estimés.

\bigskip

\subsection{Ce qu'implique l'absence de pondération \emph{a priori}}

\medskip

Renoncer à fixer des poids \emph{a priori} ne supprime pas l'hypothèse identifiée
ci-dessus : cela la \emph{déplace} d'un jugement de valeur explicite vers une propriété
statistique de l'échantillon, qui joue alors le même rôle. Chaque famille de méthodes
correspond à un postulat différent sur ce qui, dans les données, fait l'information.

\medskip

\begin{avertissement}[Le postulat implicite des pondérations endogènes]
  \label{av:cout}
  \begin{itemize}[nosep,leftmargin=1.4em]
    \item \textbf{Entropie, CRITIC, analyse en composantes principales} :
    \emph{l'information utile est la dispersion}. Une métrique qui discrimine peu entre
    les unités reçoit un poids faible, même lorsqu'elle décrit le mécanisme le plus
    important. Un risque auquel toutes les unités sont également exposées est
    mécaniquement neutralisé.
    \item \textbf{Analyse en composantes principales, axe principal} : \emph{la quantité
    d'intérêt est la direction de plus grande dispersion}. Rien ne garantit que le premier
    axe factoriel corresponde au phénomène étudié plutôt qu'à un effet de taille ou à un
    artefact d'échelle.
    \item \textbf{\emph{Benefit of the doubt}} : \emph{chaque unité est jugée sous la
    pondération qui lui est la plus favorable}. Le classement obtenu est une lecture de
    précaution, et non une mesure du niveau moyen.
    \item \textbf{Transport optimal} : \emph{la position dans le nuage résume
    l'information}, au travers d'une notion de centralité multivariée qui absorbe les
    corrélations et les échelles.
    \item \textbf{Exploration de l'espace des poids} : \emph{l'ignorance sur les poids se
    représente par une loi uniforme sur le simplexe}, ce qui est une hypothèse de plein
    droit et non une absence d'hypothèse.
  \end{itemize}
\end{avertissement}

\bigskip

Ce constat oriente la démarche suivie ici. Plutôt que de sélectionner une méthode unique,
la note décrit un ensemble de méthodes reposant sur des postulats distincts, énonce pour
chacune les propriétés vérifiées et les propriétés non vérifiées, et présente au
chapitre~\ref{sec:comparer} les procédures permettant de mesurer l'accord entre les
classements obtenus. Un désaccord marqué entre deux méthodes n'est pas un défaut à
corriger : c'est l'information selon laquelle le classement dépend du postulat retenu.

\bigskip

\subsection{Taxonomie des familles de méthodes}

\medskip

Les méthodes décrites se répartissent en six familles, selon la nature des coefficients
qu'elles introduisent et la nature de leur sortie.

\medskip

\begin{figure}[htbp]
  \centering
  \begin{tikzpicture}[
      font=\small,
      famille/.style={draw=Bleu,fill=Bleu!6,rounded corners=2pt,thick,
                      text width=3.5cm,align=center,inner sep=5pt,minimum height=1.25cm},
      racine/.style={draw=Gris,fill=GrisClair,rounded corners=2pt,thick,
                     text width=5.4cm,align=center,inner sep=5pt},
      lien/.style={draw=Gris,thick,-{Stealth}}
    ]
    \node[racine] (root) at (0,0)
      {Ordonner $\{x_1,\dots,x_n\} \subset \R^d$\\[2pt]
       \scriptsize sans hiérarchie \emph{a priori} des coordonnées};

    \node[famille] (ordres)  at (-5.6,-2.3) {\textbf{Ordres partiels}\\[2pt]
      \scriptsize dominance, fronts, comptage\\ \scriptsize\textcolor{Gris}{aucun coefficient}};
    \node[famille] (poids)   at ( 0.0,-2.3) {\textbf{Pondérations endogènes}\\[2pt]
      \scriptsize entropie, CRITIC, ACP, BoD\\ \scriptsize\textcolor{Gris}{$w$ estimé sur $X$}};
    \node[famille] (agreg)   at ( 5.6,-2.3) {\textbf{Agrégations}\\[2pt]
      \scriptsize somme, géométrique, TOPSIS\\ \scriptsize\textcolor{Gris}{compensation réglable}};

    \node[famille] (ot)      at (-5.6,-5.2) {\textbf{Rangs multivariés}\\[2pt]
      \scriptsize transport optimal\\ \scriptsize\textcolor{Gris}{position dans le nuage}};
    \node[famille] (smaa)    at ( 0.0,-5.2) {\textbf{Exploration des poids}\\[2pt]
      \scriptsize SMAA, quantile de cône\\ \scriptsize\textcolor{Gris}{loi sur $\simplex$}};
    \node[famille] (controle)at ( 5.6,-5.2) {\textbf{Comparaison et contrôle}\\[2pt]
      \scriptsize accord, stabilité, consensus\\ \scriptsize\textcolor{Gris}{pas de classement propre}};

    \draw[lien] (root) -- (ordres);
    \draw[lien] (root) -- (poids);
    \draw[lien] (root) -- (agreg);
    \draw[lien] (root) -- (ot);
    \draw[lien] (root) -- (smaa);
    \draw[lien] (root) -- (controle);

    \draw[Gris,thick,dashed,-{Stealth}] (poids) -- (agreg);
    \node[Gris,font=\scriptsize] at (2.8,-1.95) {$w$};
  \end{tikzpicture}
  \caption{Taxonomie des familles de méthodes. Les pondérations endogènes produisent un
  vecteur de poids qui alimente une fonction d'agrégation : ces deux familles se composent,
  alors que les ordres partiels et les rangs multivariés produisent directement un ordre.
  La dernière famille ne produit pas de classement mais évalue ceux des autres.}
  \label{fig:taxonomie}
\end{figure}

\bigskip

Les deux premières familles se distinguent par le nombre d'hypothèses introduites. Les
ordres partiels n'en introduisent aucune, au prix d'une incomparabilité assumée : deux
unités dont les métriques ne concordent pas ne sont pas départagées. Les pondérations
endogènes produisent un classement complet, au prix du postulat statistique correspondant.
Les rangs multivariés occupent une position intermédiaire : ils ne fixent pas de poids,
mais l'orientation du score, qui désigne la direction dans laquelle un point est considéré
comme extrême, reste une convention.

\bigskip

\subsection{Plan}

\medskip

Le chapitre~\ref{sec:cadre} fixe le cadre et les notations, et définit les propriétés qui
serviront de grille de lecture : dominance, monotonie, compensation, invariances,
conventions de rang. Le chapitre~\ref{sec:pretraitement} décrit le prétraitement
(orientation, traitement des queues de distribution, normalisation, diagnostics de
corrélation), étape dont dépendent la plupart des méthodes ultérieures.

\medskip

Les chapitres~\ref{sec:ordres} à~\ref{sec:smaa} présentent les méthodes :
ordres partiels et dominance (chapitre~\ref{sec:ordres}), pondérations endogènes
(chapitre~\ref{sec:poids}), fonctions d'agrégation (chapitre~\ref{sec:agregation}), rangs
multivariés par transport optimal (chapitre~\ref{sec:ot}) et incertitude sur les poids
(chapitre~\ref{sec:smaa}).

\medskip

Le chapitre~\ref{sec:comparer} décrit les procédures de comparaison entre classements, de
mesure de la stabilité et de construction d'un consensus. Le
chapitre~\ref{sec:application} illustre l'ensemble sur un cas d'application. Les annexes
rassemblent les notations (annexe~\ref{app:notations}), les algorithmes
(annexe~\ref{app:algos}) et la correspondance avec une implémentation de référence
(annexe~\ref{app:code}).

%=======================================================================
\section{Cadre et notations}
\label{sec:cadre}
%=======================================================================

\medskip

\subsection{Unités, métriques et groupes}

\medskip

Les objets classés sont appelés \emph{unités} ; les variables qui les décrivent sont
appelées \emph{métriques}. Aucune hypothèse n'est faite sur la nature des unités : la
seule exigence est que les $n$ unités considérées simultanément soient
\emph{comparables}, au sens où les valeurs d'une même métrique y sont mesurées sur la même
échelle et sont interprétables les unes par rapport aux autres.

\medskip

\begin{definition}[Groupe]
  \label{def:groupe}
  Un \emph{groupe} est un ensemble d'unités à l'intérieur duquel l'ordre est construit.
  Toutes les quantités estimées à partir des données --- bornes de normalisation, poids,
  moments, cartes de transport, pôles idéal et anti-idéal --- sont calculées à
  l'intérieur d'un groupe et ne sont pas transportables d'un groupe à l'autre.
\end{definition}

\bigskip

La notion de groupe est structurante. Une pondération endogène est une propriété de
l'échantillon sur lequel elle est estimée : les poids obtenus sur deux groupes différents
diffèrent, et les scores correspondants ne sont pas comparables entre groupes. Lorsqu'une
comparaison entre groupes est recherchée, elle porte sur les rangs relatifs à l'intérieur
de chaque groupe, ou sur un score estimé une seule fois sur la réunion des groupes.

\bigskip

\subsection{Notations}

\medskip

\begin{table}[htbp]
  \centering
  \small
  \renewcommand{\arraystretch}{1.2}
  \begin{tabularx}{\textwidth}{@{}l X@{}}
    \toprule
    \textbf{Symbole} & \textbf{Définition} \\
    \midrule
    $n$, $d$ & nombre d'unités (lignes) et de métriques (colonnes) d'un groupe \\
    $[n] = \{1,\dots,n\}$ & indices des unités \\
    $X = (x_{ij}) \in \R^{n\times d}$ & matrice brute ; $x_i \in \R^d$ la ligne $i$ ;
      $x_{\cdot j}$ la colonne $j$ \\
    $p \in \{-1,+1\}^d$ & polarités ; $\Xor = X\,\mathrm{diag}(p)$ la matrice orientée \\
    $\tilde{X}$ & matrice orientée puis normalisée par le schéma $\nu$ \\
    $Z$ & matrice orientée puis centrée-réduite \\
    $x_i \pareto x_k$ & dominance de Pareto stricte sur $\Xor$
      (définition~\ref{def:pareto}) \\
    $\mathcal{F}_1$, $\mathcal{F}_\ell$, $\mathcal{F}_\varepsilon$ & front, couche $\ell$,
      front $\varepsilon$ (sous-ensemble de $\mathcal{F}_1$) \\
    $\delta(i)$, $\bar\delta(i)$ & comptage de dominance ; profondeur normalisée
      $\delta/(n-1)$ \\
    $\simplex$ & simplexe des poids
      $\{w \in \R_+^d : \sum_j w_j = 1\}$ \\
    $w$, $w^{(o)}$ & poids partagés ; poids individuels de l'unité $o$ \\
    $s : \R^d \to \R$ & fonction de score, valeur élevée $=$ tête de classement \\
    $r_i$ & rang de l'unité $i$, $r = 1$ pour la valeur la plus élevée, ex \ae quo moyennés \\
    $\rho^S_{jk}$, $r_{jk}$ & corrélation de Spearman ; corrélation de Pearson \\
    $R$, $\Sigma$, $\mu$ & matrice de corrélation ; covariance (robuste) ; centre (robuste) \\
    $\lambda_m$, $a_m$, $\ell_{jm}$, $\ell'_{jm}$, $V_m$ & valeurs et vecteurs propres ;
      saturations ; saturations après rotation ; variance d'un facteur après rotation \\
    $x^+$, $x^-$ & pôles idéal et anti-idéal (quantiles marginaux $q^+$, $q^-$) \\
    $\Unif_d$, $\Bball_d$, $\Sphere^{d-1}$ & loi sphérique uniforme ; boule unité ;
      sphère unité \\
    $\Tmap$, $\bar{\Tmap}_\varepsilon$ & carte \emph{center-outward} ; son estimateur
      entropique hors échantillon \\
    $\mathcal{R}(z) = \norm{\Tmap(z)}$, $\mathcal{S}(z) = \Tmap(z)/\norm{\Tmap(z)}$ &
      rang et signe \emph{center-outward} \\
    $u^*$, $\theta_0$, $\alpha$, $\hat{r}_\alpha$, $\hat{q}_{1-\alpha}$ & direction
      d'intérêt ; demi-angle du cône ; niveau ; seuil radial ; seuil projeté \\
    $\varepsilon$, $m$ & régularisation entropique ; taille de la grille de référence \\
    $\tau_b$, $\tau_w$, $W$, $\mathrm{RBO}_p$, $J_k$ & Kendall ; Kendall pondéré ;
      concordance ; \emph{rank-biased overlap} ; recouvrement des $k$ premiers \\
    $V_{\mathrm{strict}}(s)$, $V_{\mathrm{tie}}(s)$ & taux d'inversions strictes ; taux
      d'ex \ae quo sur les paires comparables \\
    $B$, $T$ & nombre de tirages bootstrap ; nombre de tirages de poids \\
    $\rho$ & borne des rapports de poids d'une région d'assurance \\
    $\kappa$, $\delta$ & paramètre de parts ; plancher additif (option) \\
    \bottomrule
  \end{tabularx}
  \caption{Notations employées dans l'ensemble de la note. Les symboles propres à une
  méthode sont redéfinis au chapitre correspondant.}
  \label{tab:notations}
\end{table}

\bigskip

\subsection{Polarité et orientation}

\medskip

\begin{definition}[Polarité]
  \label{def:polarite}
  Une métrique $j$ est de \emph{polarité positive} lorsqu'une valeur élevée place l'unité
  en tête du classement recherché, et de \emph{polarité négative} dans le cas contraire.
  Le vecteur $p \in \{-1,+1\}^d$ rassemble les polarités.
\end{definition}

\bigskip

L'ensemble des méthodes décrites suppose que toutes les métriques sont de polarité
positive. L'\emph{orientation} est l'opération qui ramène la matrice brute à cette
convention : $\Xor = X\,\mathrm{diag}(p)$, soit $\mathring{x}_{ij} = p_j\,x_{ij}$. Cette
transformation est affine ; la section~\ref{sec:orientation} compare les variantes
possibles et justifie ce choix.

\bigskip

\subsection{Dominance de Pareto}

\medskip

\begin{definition}[Dominance de Pareto]
  \label{def:pareto}
  Sur la matrice orientée $\Xor$, l'unité $i$ \emph{domine strictement} l'unité $k$, ce
  qui se note $x_i \pareto x_k$, lorsque

  \medskip

  \[
    \forall j \in [d],\ \mathring{x}_{ij} \ge \mathring{x}_{kj}
    \quad \text{et} \quad
    \exists j_0 \in [d],\ \mathring{x}_{ij_0} > \mathring{x}_{kj_0}.
  \]

  \bigskip

  La \emph{dominance faible}, notée $x_i \paretoeq x_k$, retient la seule première
  condition.
\end{definition}

\bigskip

La relation $\pareto$ est un ordre strict partiel : elle est irréflexive, transitive et
antisymétrique, mais non totale dès que $d \ge 2$. Deux unités dont les métriques ne
concordent pas sont \emph{incomparables} au sens de $\pareto$. Géométriquement, $i$ domine
$k$ si et seulement si $\mathring{x}_i$ appartient au cône $\mathring{x}_k + \R_+^d$ privé
de son sommet.

\medskip

\begin{figure}[htbp]
  \centering
  \begin{tikzpicture}[scale=1.0]
    \draw[->,thick,Gris] (0,0) -- (7.2,0) node[right,black,font=\small] {métrique $1$};
    \draw[->,thick,Gris] (0,0) -- (0,5.0) node[above,black,font=\small] {métrique $2$};

    % cône supérieur : points dominant x_i
    \begin{scope}[opacity=0.16]
      \fill[Bleu] (3.5,2.2) rectangle (7.0,4.8);
    \end{scope}
    % cône inférieur : points dominés par x_i
    \begin{scope}[opacity=0.16]
      \fill[Vert] (0.05,0.05) rectangle (3.5,2.2);
    \end{scope}

    \foreach \x/\y in {0.8/1.0, 1.5/0.6, 2.0/1.9, 2.9/1.2, 3.1/0.5, 2.6/2.2}
      \fill[Vert] (\x,\y) circle (2.1pt);
    \foreach \x/\y in {4.2/2.6, 5.4/3.1, 4.0/2.9, 6.1/2.4}
      \fill[Bleu] (\x,\y) circle (2.1pt);
    \foreach \x/\y in {1.2/2.4, 0.6/3.0, 2.4/3.0, 4.9/1.8, 5.9/0.9, 1.0/4.2}
      \fill[Gris] (\x,\y) circle (2.1pt);

    \draw[Bleu,thick,-{Stealth}] (3.5,2.2) -- (4.8,2.2);
    \draw[Bleu,thick,-{Stealth}] (3.5,2.2) -- (3.5,3.5);
    \fill[Rouge] (3.5,2.2) circle (2.8pt) node[below right,font=\scriptsize,black] {$x_i$};

    \node[Bleu,font=\scriptsize,align=center] at (5.9,4.35)
      {unités dominant $x_i$};
    \node[Vert,font=\scriptsize,align=center] at (1.55,0.25)
      {unités dominées par $x_i$};
    \node[Gris,font=\scriptsize,align=center] at (1.35,3.55)
      {incomparables};
  \end{tikzpicture}
  \caption{Cône de dominance en dimension $2$, les deux métriques étant orientées en
  polarité positive. Le plan se partage en trois régions : le cône supérieur droit issu
  de $x_i$ (unités qui le dominent), le cône inférieur gauche (unités qu'il domine) et les
  deux quadrants restants (unités incomparables avec $x_i$). La proportion de la troisième
  région croît rapidement avec $d$.}
  \label{fig:cone}
\end{figure}

\bigskip

\subsection{Monotonie stricte et monotonie faible}

\medskip

La compatibilité d'un score avec la dominance se décline en deux propriétés distinctes,
qu'il est nécessaire de séparer : plusieurs méthodes vérifient la seconde sans vérifier la
première, et la différence porte précisément sur le traitement des ex \ae quo.

\medskip

\begin{definition}[Monotonie stricte et monotonie faible]
  \label{def:monotonie}
  Une fonction de score $s : \R^d \to \R$ est

  \medskip

  \begin{itemize}[nosep,leftmargin=1.4em]
    \item \emph{strictement monotone} pour $\pareto$ lorsque
    $x_i \pareto x_k \Rightarrow s(x_i) > s(x_k)$ ;
    \item \emph{faiblement monotone} lorsque
    $x_i \pareto x_k \Rightarrow s(x_i) \ge s(x_k)$.
  \end{itemize}
\end{definition}

\bigskip

La monotonie stricte est la seule propriété exigible sans engager de jugement de valeur :
elle formalise l'idée qu'une unité meilleure sur une métrique et non pire sur les autres
ne peut pas être classée plus bas. Une somme pondérée à poids strictement positifs est
strictement monotone ; la même somme avec un poids nul ne l'est que faiblement, puisque
deux unités ne différant que sur la métrique de poids nul reçoivent le même score.

\medskip

\begin{remarque}[Conséquence pour les contrôles]
  \label{rem:monotonie-controle}
  La distinction entre les deux notions se traduit dans le contrôle de cohérence du
  chapitre~\ref{sec:comparer} par deux quantités séparées : le taux d'inversions strictes
  $V_{\mathrm{strict}}(s)$, qui compte les paires dominantes pour lesquelles
  $s(x_i) < s(x_k)$, et le taux d'ex \ae quo $V_{\mathrm{tie}}(s)$, qui compte celles pour
  lesquelles $s(x_i) = s(x_k)$. Une méthode faiblement monotone a
  $V_{\mathrm{strict}} = 0$ et $V_{\mathrm{tie}} \ge 0$ ; leur confusion en un indicateur
  unique rendrait incomparables des méthodes qui ne le sont pas pour la même raison.
\end{remarque}

\bigskip

\subsection{Compensation}

\medskip

\begin{definition}[Compensation]
  \label{def:compensation}
  Une fonction d'agrégation $s$ est \emph{compensatoire} lorsqu'une dégradation sur une
  coordonnée peut être intégralement compensée par une amélioration sur une autre,
  c'est-à-dire lorsque, pour tout $x$, toute paire $(j,k)$ et toute perturbation
  $-h$ sur la coordonnée $j$, il existe une perturbation $h'$ sur la coordonnée $k$
  laissant $s$ inchangé.
\end{definition}

\bigskip

La compensation est une propriété graduelle plutôt que binaire ; les exemples suivants en
délimitent les extrêmes.

\medskip

\begin{itemize}[nosep,leftmargin=1.4em]
  \item La somme pondérée $s(x) = \sum_j w_j x_j$ est totalement compensatoire : le taux
  de substitution entre les métriques $j$ et $k$ vaut $w_k/w_j$, constant sur tout
  l'espace.
  \item L'ordre lexicographique n'est pas compensatoire : aucune amélioration sur la
  deuxième coordonnée ne rattrape une dégradation sur la première.
  \item La moyenne géométrique $s(x) = \prod_j x_j^{w_j}$ est partiellement
  compensatoire : le taux de substitution dépend du point, et une coordonnée proche de
  zéro ne peut plus être compensée.
  \item Les indices à pénalité explicite, qui retranchent une fonction croissante de la
  dispersion des coordonnées, et les méthodes fondées sur le regret maximal occupent des
  positions intermédiaires paramétrables.
\end{itemize}

\bigskip

Le degré de compensation admissible est un choix méthodologique et non une propriété des
données. Il est explicité pour chaque méthode au chapitre correspondant, et fait partie
des dimensions de comparaison du chapitre~\ref{sec:comparer}.

\bigskip

\subsection{Invariances}

\medskip

Trois familles de transformations servent à caractériser les méthodes. Une méthode est
dite \emph{invariante} pour une transformation lorsque l'ordre qu'elle produit est
inchangé lorsqu'on l'applique aux données.

\medskip

\begin{definition}[Invariances]
  \label{def:invariances}
  Soit $\mathcal{A}$ une méthode associant à $X$ un ordre sur $[n]$.

  \medskip

  \begin{itemize}[nosep,leftmargin=1.4em]
    \item $\mathcal{A}$ est \emph{invariante affine par coordonnée} lorsque l'ordre est
    inchangé par $x_{ij} \mapsto a_j x_{ij} + b_j$ pour tout $a_j > 0$ et $b_j \in \R$.
    \item $\mathcal{A}$ est \emph{invariante monotone par coordonnée} lorsque l'ordre est
    inchangé par $x_{ij} \mapsto g_j(x_{ij})$ pour toute fonction $g_j$ strictement
    croissante. Cette invariance implique la précédente.
    \item $\mathcal{A}$ est \emph{invariante par permutation des coordonnées} lorsque
    l'ordre est inchangé lorsqu'on renumérote les métriques, c'est-à-dire lorsque la
    méthode ne privilégie aucune métrique par construction.
  \end{itemize}
\end{definition}

\bigskip

L'invariance par permutation est vérifiée par toutes les méthodes considérées ici, ce qui
est la traduction formelle de l'absence de hiérarchie \emph{a priori}. Les deux autres
invariances les séparent nettement : la dominance de Pareto et tout ce qui s'en déduit
sont invariants monotones ; les méthodes fondées sur la covariance sont invariantes
affines mais pas monotones ; les méthodes fondées sur la normalisation min--max ne sont
en général ni l'une ni l'autre, puisqu'elles dépendent des valeurs extrêmes observées.

\bigskip

\subsection{Rangs et ex \ae quo}

\medskip

\begin{definition}[Rang]
  \label{def:rang}
  Étant donné un score $s$ et un groupe de $n$ unités, le \emph{rang} $r_i$ de l'unité $i$
  est sa position dans le tri décroissant de $(s(x_1),\dots,s(x_n))$ : la valeur la plus
  élevée reçoit le rang $1$. Les unités de score égal reçoivent la \emph{moyenne} des
  positions qu'elles occupent, de sorte que $\sum_i r_i = n(n+1)/2$ quelle que soit la
  structure des ex \ae quo.
\end{definition}

\bigskip

Cette convention, dite des rangs moyens, est celle qu'emploient les corrélations de
Spearman et de Kendall corrigées des ex \ae quo. Elle est employée dans toute la note, y
compris dans les normalisations par rangs de la section~\ref{sec:normalisation}, où les
rangs sont pris dans l'ordre \emph{croissant} de la métrique orientée afin qu'une valeur
élevée corresponde à une valeur normalisée élevée.

\medskip

\begin{remarque}[Ex \ae quo légitimes]
  \label{rem:exaequo}
  Les ex \ae quo ne sont pas tous des artefacts. Certaines méthodes en produisent par
  construction : les scores à valeurs discrètes, les scores bornés atteints par plusieurs
  unités, ou les rangs de couche du tri non dominé. D'autres en produisent par effet de
  bord du prétraitement, notamment la winsorisation
  (section~\ref{sec:queues}). La distinction est conservée dans tous les diagnostics.
\end{remarque}

%=======================================================================
\section{Prétraitement}
\label{sec:pretraitement}
%=======================================================================

\medskip

Le prétraitement transforme la matrice brute $X$ en une matrice $\tilde{X}$ utilisable par
les méthodes des chapitres suivants. Il comporte quatre opérations : l'orientation, le
traitement éventuel des queues de distribution, la normalisation, et le calcul des
diagnostics qui documentent la structure de corrélation. Ce n'est pas une étape neutre :
le choix de normalisation modifie le classement final au même titre que le choix de la
méthode d'agrégation, et relève à ce titre de l'analyse de sensibilité.

\bigskip

\subsection{Orientation}
\label{sec:orientation}

\medskip

L'orientation ramène toutes les métriques en polarité positive
(définition~\ref{def:polarite}). Pour une métrique de polarité négative, trois
transformations sont couramment employées :

\medskip

\begin{itemize}[nosep,leftmargin=1.4em]
  \item la \emph{négation} $x_{ij} \mapsto -x_{ij}$, affine et définie sans condition ;
  \item le \emph{complément à la borne supérieure}
  $x_{ij} \mapsto \max_k x_{kj} - x_{ij}$, affine également, mais dépendant d'une
  statistique de l'échantillon, donc non transportable d'un groupe à l'autre ;
  \item l'\emph{inversion} $x_{ij} \mapsto 1/x_{ij}$, définie pour une métrique
  strictement positive, strictement décroissante mais non affine.
\end{itemize}

\bigskip

Ces trois transformations induisent le même ordre marginal sur la métrique $j$, donc la
même relation de dominance ; elles ne sont en revanche pas équivalentes pour les méthodes
qui exploitent la structure de covariance. Seules les deux premières sont affines, et donc
seules elles laissent invariants les résultats de l'analyse en composantes principales, de
la distance de Mahalanobis et du transport optimal dans le cas elliptique. La négation est
la seule des trois à ne dépendre d'aucune statistique de l'échantillon : c'est la variante
retenue par défaut dans la suite, l'orientation étant alors $\Xor = X\,\mathrm{diag}(p)$.

\bigskip

\subsection{Queues de distribution}
\label{sec:queues}

\medskip

Les métriques de concentration produisent fréquemment des distributions très
asymétriques : une poignée d'unités atteignent des valeurs extrêmes tandis que la masse se
concentre près de la borne inférieure. Trois attitudes sont possibles, chacune ayant un
coût identifiable.

\medskip

\paragraph{Aucun traitement.} Les valeurs extrêmes sont conservées telles quelles. C'est
le choix cohérent avec l'objectif lorsque les unités extrêmes sont précisément celles que
le classement doit identifier. Son coût est que la normalisation min--max est alors pilotée
par une seule unité par métrique, et que les méthodes fondées sur la covariance sont
sensibles à quelques points de levier.

\bigskip

\paragraph{Winsorisation.} La transformation
$x_{ij} \mapsto \min\!\big(x_{ij},\,\hat{F}_j^{-1}(q)\big)$ rabat les valeurs supérieures
au quantile empirique d'ordre $q$ sur ce quantile. Elle stabilise les statistiques
d'échelle et les estimateurs de covariance, mais elle a un coût direct sur la
discrimination dans la queue haute, qui est l'objet même du classement.

\medskip

\begin{remarque}[Coût de la winsorisation en ex \ae quo]
  \label{rem:winsor}
  Après winsorisation au quantile $q$ puis normalisation min--max, les
  $\lceil (1-q)\,n \rceil$ unités les plus élevées de chaque métrique reçoivent exactement
  la valeur $1$. Pour $n = 5\,000$ et $q = 0{,}99$, cinquante unités sont ex \ae quo au
  maximum de chaque métrique et ne sont plus départagées, métrique par métrique, par aucune
  méthode compensatoire. La discrimination perdue l'est définitivement : aucune étape
  ultérieure ne la restaure. La winsorisation est donc désactivée par défaut ($q$ non
  renseigné) ; lorsqu'elle est activée, la valeur de $q$ est rapportée avec le nombre
  d'ex \ae quo induits.
\end{remarque}

\bigskip

\paragraph{Transformation monotone.} Une transformation strictement croissante appliquée
coordonnée par coordonnée --- $\log(1+x)$, racine carrée, transformation de Box--Cox dont
le paramètre est estimé par maximum de vraisemblance --- réduit l'asymétrie sans créer
d'ex \ae quo. Elle est adaptée lorsque l'objectif est de stabiliser les estimateurs de
covariance tout en conservant l'ordre marginal complet. Son coût est une modification des
écarts cardinaux, donc des taux de substitution implicites d'une somme pondérée.

\bigskip

\begin{remarque}[Invariance de la dominance]
  \label{rem:invariance-front}
  Une transformation strictement croissante appliquée coordonnée par coordonnée laisse la
  relation $\pareto$ inchangée : la dominance ne dépend que des ordres marginaux. Le front
  de Pareto, les couches et le comptage de dominance
  (chapitre~\ref{sec:ordres}) sont donc insensibles au choix entre « aucun traitement » et
  « transformation monotone », ainsi qu'au choix du schéma de normalisation. La
  winsorisation fait exception : elle n'est pas strictement croissante, elle crée des
  ex \ae quo, et elle modifie donc la relation de dominance elle-même.
\end{remarque}

\bigskip

\subsection{Schémas de normalisation}
\label{sec:normalisation}

\medskip

La normalisation ramène les $d$ métriques sur une échelle commune. Elle est nécessaire dès
lors qu'une méthode combine les coordonnées, et le schéma retenu détermine les invariances
disponibles ainsi que les méthodes admissibles en aval.

\medskip

\begin{table}[htbp]
  \centering
  \scriptsize
  \renewcommand{\arraystretch}{1.35}
  \begin{tabularx}{\textwidth}{@{}l l Y l Y@{}}
    \toprule
    \textbf{Schéma} & \textbf{Formule} & \textbf{Effet} & \textbf{Invariances} &
    \textbf{Méthodes compatibles} \\
    \midrule
    Min--max &
    $\tilde{x}_{ij} = \dfrac{\mathring{x}_{ij} - \min_k \mathring{x}_{kj}}
                            {\max_k \mathring{x}_{kj} - \min_k \mathring{x}_{kj}}$ &
    Borne dans $[0,1]$ et conserve les écarts relatifs ; une seule unité atypique comprime
    toute la métrique &
    affine &
    entropie, CRITIC, \emph{benefit of the doubt}, moyenne géométrique, indices à
    pénalité, TOPSIS, VIKOR \\
    Centrage--réduction &
    $\tilde{x}_{ij} = \dfrac{\mathring{x}_{ij} - \bar{x}_j}{\sigma_j}$ &
    Égalise les variances, donc égalise implicitement les poids d'une somme ; non borné,
    conserve l'asymétrie &
    affine &
    ACP, Mahalanobis, projection blanchie, transport optimal \\
    Robuste &
    $\tilde{x}_{ij} = \dfrac{\mathring{x}_{ij} - \med_j}{\MAD_j}$ &
    Comme la précédente, mais insensible aux valeurs extrêmes ; adaptée aux distributions
    très asymétriques &
    affine &
    ACP, Mahalanobis, projection blanchie, transport optimal \\
    Rangs &
    $\tilde{x}_{ij} = \dfrac{\rg(\mathring{x}_{ij}) - 1}{n-1}$ &
    Détruit l'information cardinale : un écart important et un écart infime deviennent
    identiques ; borne dans $[0,1]$ &
    monotone &
    toutes ; seul schéma rendant le classement invariant par transformation monotone \\
    Quantile gaussien &
    $\tilde{x}_{ij} = \Phi^{-1}\!\Big(\dfrac{\rg(\mathring{x}_{ij}) - 1/2}{n}\Big)$ &
    Marges gaussiennes, structure de dépendance (copule) préservée ; non borné &
    monotone &
    ACP, Mahalanobis, transport optimal \\
    \bottomrule
  \end{tabularx}
  \caption{Schémas de normalisation. Les rangs sont pris dans l'ordre croissant de la
  métrique orientée, ex \ae quo moyennés. Aucun schéma n'est neutre : le
  centrage--réduction constitue déjà une pondération implicite, puisqu'il égalise la
  contribution de chaque métrique à la variance d'une somme non pondérée.}
  \label{tab:normalisation}
\end{table}

\bigskip

Deux contraintes fixent le schéma admissible pour une méthode donnée. La première est le
\emph{signe} : les méthodes exigeant des données positives --- moyenne géométrique,
indices à pénalité, \emph{benefit of the doubt} --- n'admettent que les schémas bornés
dans $[0,1]$. La seconde est la \emph{translation} : les pondérations par entropie
dépendent de l'origine de l'échelle, ce qui exclut les schémas centrés. À l'inverse, les
méthodes fondées sur la covariance perdent leur interprétation sur des données min--max,
dont les variances marginales sont pilotées par l'étendue observée.

\bigskip

\subsection{Diagnostics de corrélation}
\label{sec:diagnostics}

\medskip

La structure de corrélation des métriques normalisées conditionne l'adéquation des
familles de méthodes. Quatre diagnostics la résument ; ils sont calculés sur la matrice de
corrélation de Spearman $\rho^S = (\rho^S_{jk})$, qui ne dépend que des ordres marginaux
et reste donc stable sous transformation monotone et en présence de valeurs extrêmes.

\medskip

\paragraph{Corrélation moyenne.} La statistique
$\bar{\rho} = \tfrac{2}{d(d-1)} \sum_{j<k} \abs{\rho^S_{jk}}$ résume l'intensité de la
redondance entre métriques. Une valeur proche de $0$ indique des métriques quasi
indépendantes ; une valeur élevée indique que les métriques mesurent en grande partie la
même chose, et qu'une somme pondérée compte alors plusieurs fois la même information.

\bigskip

\paragraph{Indice KMO.} L'indice de \citet{kaiser1974} compare les corrélations simples
aux corrélations partielles :

\medskip

\[
  \mathrm{KMO} =
  \frac{\sum_{j \ne k} r_{jk}^2}{\sum_{j \ne k} r_{jk}^2 + \sum_{j \ne k} \pi_{jk}^2},
\]

\bigskip

où $\pi_{jk}$ est la corrélation partielle entre les métriques $j$ et $k$, les autres
étant contrôlées. Une valeur élevée signale que les corrélations observées s'expliquent
par un petit nombre de facteurs communs, et donc qu'une pondération de type factoriel est
identifiable ; une valeur faible signale que les corrélations sont essentiellement des
liens deux à deux, situation dans laquelle un modèle factoriel n'apporte rien.

\bigskip

\paragraph{Test de sphéricité de Bartlett.} Le test de \citet{bartlett1951} compare la
matrice de corrélation à l'identité, au moyen de la statistique

\medskip

\[
  \chi^2 = -\Big(n - 1 - \tfrac{2d+5}{6}\Big)\,\ln \det R,
  \qquad \text{de loi } \chi^2 \text{ à } \tfrac{d(d-1)}{2} \text{ degrés de liberté.}
\]

\bigskip

Le non-rejet indique une matrice de corrélation compatible avec l'indépendance, donc
l'absence de structure factorielle à extraire. Sur les grands échantillons, le test
rejette presque systématiquement ; sa lecture est alors principalement qualitative et
l'indice KMO est plus informatif.

\bigskip

\paragraph{Lecture conjointe.} Les trois configurations suivantes appellent des familles de
méthodes différentes.

\medskip

\begin{itemize}[nosep,leftmargin=1.4em]
  \item \textbf{Métriques fortement corrélées} ($\bar{\rho}$ élevé, KMO élevé). Une somme
  pondérée surpondère de fait le mécanisme commun. Les méthodes qui décorrèlent
  explicitement --- analyse en composantes principales, distance de Mahalanobis, CRITIC,
  transport optimal --- sont adaptées à cette configuration.
  \item \textbf{Métriques quasi indépendantes} ($\bar{\rho}$ faible, KMO faible).
  L'analyse en composantes principales n'a pas d'axe dominant et son premier axe est
  instable. Les pondérations par dispersion et les ordres partiels sont adaptés, mais le
  front de Pareto est alors volumineux (chapitre~\ref{sec:ordres}).
  \item \textbf{Corrélations négatives marquées.} Les métriques décrivent des mécanismes
  antagonistes, et la réduction à un score unique perd une part importante de
  l'information. Plusieurs scores thématiques, ou une méthode admettant
  l'incomparabilité, sont alors plus fidèles aux données.
\end{itemize}

\bigskip

\subsection{Algorithme de prétraitement}

\medskip

\begin{algorithm}[htbp]
  \caption{Prétraitement}
  \label{algo:pretraitement}
  \begin{algorithmic}[1]
    \Require matrice brute $X \in \R^{n \times d}$ ; polarités $p \in \{-1,+1\}^d$ ;
      quantile de winsorisation $q$ (éventuellement non renseigné) ; transformation
      monotone $g$ (éventuellement l'identité) ; schéma de normalisation $\nu$
    \Ensure matrice normalisée $\tilde{X}$ ; matrice de corrélation $\rho^S$ ;
      diagnostics $(\bar{\rho},\mathrm{KMO},\chi^2)$
    \State écarter les lignes comportant au moins une valeur manquante, et signaler leur
      nombre \Comment{Convention de lignes complètes}
    \For{$j = 1$ \textbf{à} $d$}
      \State $\mathring{x}_{\cdot j} \gets p_j \cdot x_{\cdot j}$
        \Comment{Orientation en polarité positive}
      \If{$q$ est renseigné}
        \State $\mathring{x}_{\cdot j} \gets
          \min\!\big(\mathring{x}_{\cdot j},\, \hat{F}_j^{-1}(q)\big)$
        \Comment{Winsorisation ; le nombre d'ex \ae quo créés est rapporté}
      \EndIf
      \State $\mathring{x}_{\cdot j} \gets g(\mathring{x}_{\cdot j})$
        \Comment{Transformation monotone éventuelle}
      \If{$\mathring{x}_{\cdot j}$ est constante}
        \State signaler la colonne $j$ et l'exclure des méthodes qui divisent par son
          échelle
      \EndIf
      \State $\tilde{x}_{\cdot j} \gets \nu(\mathring{x}_{\cdot j})$
        \Comment{Schéma retenu, table~\ref{tab:normalisation}}
    \EndFor
    \State $\rho^S \gets \operatorname{corr}_{\text{Spearman}}(\tilde{X})$
    \State $\bar{\rho} \gets \frac{2}{d(d-1)}\sum_{j<k} \abs{\rho^S_{jk}}$ ;
      calculer $\mathrm{KMO}$ et la statistique de Bartlett
    \State \Return $\tilde{X}$, $\rho^S$, $(\bar{\rho},\mathrm{KMO},\chi^2)$
  \end{algorithmic}
\end{algorithm}

\bigskip

L'ordre des opérations n'est pas indifférent. L'orientation précède tout le reste, afin
que les quantiles et les statistiques d'échelle soient calculés sur des colonnes déjà
mises en polarité positive. La winsorisation précède la normalisation, faute de quoi les
bornes min--max seraient fixées par des valeurs ultérieurement rabattues. Les diagnostics
sont calculés après normalisation, mais leur valeur ne dépend du schéma que pour les
statistiques fondées sur Pearson : la matrice de Spearman est invariante pour tous les
schémas de la table~\ref{tab:normalisation} à l'exception de la winsorisation.

%=======================================================================
\section{Ordres partiels et dominance}
\label{sec:ordres}
%=======================================================================

% À rédiger : prompt P13

%=======================================================================
\section{Pondérations endogènes}
\label{sec:poids}
%=======================================================================

% À rédiger : prompt P13

%=======================================================================
\section{Fonctions d'agrégation}
\label{sec:agregation}
%=======================================================================

% À rédiger : prompt P14

%=======================================================================
\section{Rangs multivariés par transport optimal}
\label{sec:ot}
%=======================================================================

% À rédiger : prompt P15

%=======================================================================
\section{Incertitude sur les poids}
\label{sec:smaa}
%=======================================================================

% À rédiger : prompt P14

%=======================================================================
\section{Comparer et contrôler}
\label{sec:comparer}
%=======================================================================

% À rédiger : prompt P16

%=======================================================================
\section{Application}
\label{sec:application}
%=======================================================================

% À rédiger : prompt P16

%=======================================================================
\appendix
%=======================================================================

%=======================================================================
\section{Notations}
\label{app:notations}
%=======================================================================

% À rédiger : prompt P16 (reprise récapitulative de la table~\ref{tab:notations})

%=======================================================================
\section{Algorithmes récapitulatifs}
\label{app:algos}
%=======================================================================

% À rédiger : prompt P16

%=======================================================================
\section{Correspondance document -- code}
\label{app:code}
%=======================================================================

% À rédiger : prompt P16

%=======================================================================
% Annexe D — bibliographie (le titre est produit par \bibsection, redéfini
% dans le préambule pour que l'annexe soit numérotée comme les autres).
%=======================================================================

\nocite{*}
\bibliographystyle{plainnat}
\bibliography{bibliography}

\end{document}
